\documentclass[12pt,a4paper]{article}

\usepackage[a4paper,margin=1in]{geometry}
\usepackage[final]{graphicx}
\graphicspath{{output_plots/}}
\usepackage{booktabs}
\usepackage{amsmath,amssymb}
\usepackage{float}
\usepackage{hyperref}
\usepackage{xcolor}
\usepackage{listings}
\usepackage{enumitem}
\usepackage{fancyhdr}
\usepackage{longtable}

\hypersetup{colorlinks=true,linkcolor=blue,urlcolor=blue}

\pagestyle{fancy}
\fancyhf{}
\lhead{ICS1512}
\rhead{Machine Learning Algorithms Laboratory}
\cfoot{\thepage}

\lstset{
language=Python,
basicstyle=\ttfamily\small,
numbers=left,
frame=single,
breaklines=true,
showstringspaces=false
}

\begin{document}

\begin{tabular}{|p{4cm}|p{11cm}|}
\hline
Experiment Title & Email Spam/Ham Classification using Na\"ive Bayes and KNN
\footnote{\textbf{GitHub Repository: }\url{https://github.com/Danush-k/Machine-Learning}}\\ \hline
Student Name & Danusu K \\ \hline
Register Number & 3122247001013 \\ \hline
Faculty & Dr. Poreddy Ajay Kumar Reddy \\ \hline
Submission Date & \today \\ \hline
\end{tabular}

\section{Objective}
Build and evaluate binary classifiers for email spam/ham detection using Na\"ive Bayes and K-Nearest Neighbors (KNN) algorithms. Compare Gaussian, Multinomial, and Bernoulli Na\"ive Bayes variants, and study the effect of varying the neighborhood parameter $k$. Perform spatial index structure comparisons (KDTree vs. BallTree), execute hyperparameter optimization (GridSearchCV vs. RandomizedSearchCV), and analyze theoretical and practical time complexities.

\section{Problem Statement}
Email spam poses significant security threats and performance issues. The problem is to develop a machine learning pipeline to classify emails as Ham (legitimate, class 0) or Spam (unsolicited bulk, class 1). The input is a feature vector of word and character occurrences from the Spambase dataset, and the output is the binary prediction. The objective is to design classifiers that achieve high predictive performance (accuracy, precision, recall, and F1-score) while minimizing prediction times.

\section{Dataset Description}
\begin{longtable}{|p{6cm}|p{8cm}|}
\hline
Dataset Name & Spambase Dataset \\ \hline
Dataset Source & UCI Machine Learning Repository \\ \hline
Number of Samples & 4601 \\ \hline
Number of Features & 57 \\ \hline
Number of Classes & 2 \\ \hline
Missing Values & 0 \\ \hline
Train-Test Split & 80\% Train, 20\% Test \\ \hline
\end{longtable}

\section{Exploratory Data Analysis}
Exploratory Data Analysis was performed to evaluate the structure, distribution, and relationship patterns within the Spambase dataset:
\begin{itemize}
\item \textbf{Dataset Overview}: Inspecting data shapes, columns, and target fields.
\item \textbf{Statistical Summary}: Computing features' ranges, mean values, standard deviations, and skewness.
\item \textbf{Missing Value Analysis}: Scanning the dataset to ensure completeness.
\item \textbf{Class Distribution}: Evaluating target variable balance.
\item \textbf{Correlation Matrix}: Analyzing linear dependency between the predictor columns and target class.
\item \textbf{Feature Distribution}: Inspecting spread and density shape of individual key features.
\end{itemize}

\begin{figure}[H]
\centering
\includegraphics[width=0.6\linewidth]{1_class_distribution.png}
\caption{Spam vs. Ham Class Count Distribution}
\label{fig:class_dist}
\end{figure}

\textbf{Inference:}
The dataset contains 4,601 instances with a target distribution of approximately 60.6\% ham (2,788 instances) and 39.4\% spam (1,813 instances). The class distribution shows that while there is a slight imbalance, it is not severe enough to require special resampling or class-weight adjustment. Performance evaluation can rely safely on conventional accuracy, precision, and F1 metrics.

\begin{figure}[H]
\centering
\includegraphics[width=0.75\linewidth]{2_correlation_heatmap.png}
\caption{Correlation Heatmap of Top Features with Class Target}
\label{fig:corr_heatmap}
\end{figure}

\textbf{Inference:}
The heatmap displays the Pearson correlation values for the 12 features most strongly associated with the target class. Features such as \texttt{word\_freq\_your}, \texttt{word\_freq\_000}, and \texttt{char\_freq\_\$} exhibit strong positive correlations with the spam class, indicating they occur frequently in spam emails. In contrast, personal identifier frequencies show weak negative correlations, making them indicators of ham emails.

\begin{figure}[H]
\centering
\includegraphics[width=0.85\linewidth]{3_feature_histograms.png}
\caption{Histograms of Most Important Predictor Features}
\label{fig:histograms}
\end{figure}

\textbf{Inference:}
The histograms compare the frequency distributions of the top 6 predictor features across the two classes. Most frequency values are close to zero for ham emails, demonstrating the sparsity of commercial keywords in regular text. Spam emails exhibit significantly wider distributions with heavier tails, suggesting that these terms occur repeatedly when a message is unsolicited.

\begin{figure}[H]
\centering
\includegraphics[width=0.85\linewidth]{4_feature_boxplots.png}
\caption{Boxplots of Key Features Split by Spam and Ham Class}
\label{fig:boxplots}
\end{figure}

\textbf{Inference:}
The boxplots represent the median, quartiles, and outlier thresholds for the top 6 predictor features. Legitimate emails have extremely narrow distributions with medians very close to zero, reflecting the absence of commercial triggers. Spam emails have higher medians and wider spreads, which allows decision models to leverage these feature thresholds during classification.

\section{Data Preprocessing}
Explain:
\begin{itemize}
\item \textbf{Missing Value Handling}: Explicit scanning revealed zero missing values in the dataset. If missing values were present, numerical fields would be filled with their column median, and categorical fields with their column mode.
\item \textbf{Encoding}: Categorical fields are converted using a label encoder. All predictor features in the Spambase dataset are continuous frequency rates, so only target label encoding validation was required.
\item \textbf{Feature Scaling}: Since distance calculations in KNN are sensitive to feature scales, min-max normalization is applied to scale all features into the $[0,1]$ range:
$$x_{\text{scaled}} = \frac{x - x_{\text{min}}}{x_{\text{max}} - x_{\text{min}}}$$
\item \textbf{Train-Test Split}: The preprocessed dataset was split into 80\% training data (3,680 samples) and 20\% test data (921 samples) using a random seed of 42 to ensure reproducibility.
\item \textbf{Feature Engineering}: Features are scaled directly based on their original occurrence frequency counts to preserve relative density.
\end{itemize}

\section{Machine Learning Algorithm}
Explain:
\begin{itemize}
\item \textbf{Working Principle}: 
Na\"ive Bayes is a generative probabilistic classifier that computes the class-conditional probabilities of features using Bayes' Theorem under the assumption that features are conditionally independent given the class label. KNN is an instance-based, non-parametric lazy learner that classifies query points by finding the $k$ nearest neighbors in the feature space and performing a majority vote.
\item \textbf{Mathematical Formulation}:
For Na\"ive Bayes, the posterior probability of class $C$ given feature vector $\mathbf{x}$ is:
$$P(C | \mathbf{x}) = \frac{P(\mathbf{x} | C) P(C)}{P(\mathbf{x})} \propto P(C) \prod_{i=1}^{d} P(x_i | C)$$
Three variants are used: Gaussian NB ($x_i|C \sim \mathcal{N}(\mu, \sigma^2)$), Multinomial NB ($P(x_i|C) \propto \theta_i^{x_i}$), and Bernoulli NB ($P(x_i|C) = p_i^{x_i}(1-p_i)^{1-x_i}$).

For KNN, distance is computed using Minkowski metric:
$$D(\mathbf{x}, \mathbf{y}) = \left( \sum_{j=1}^{d} |x_j - y_j|^p \right)^{1/p}$$
where $p=1$ represents Manhattan distance and $p=2$ represents Euclidean distance. The decision is:
$$\hat{y} = \arg\max_{c} \sum_{\mathbf{x}_i \in \mathcal{N}_k(\mathbf{x}_q)} w_i I(y_i = c)$$
where $w_i = 1$ for uniform voting, and $w_i = \frac{1}{D(\mathbf{x}_q, \mathbf{x}_i) + \epsilon}$ for distance-weighted voting. Tree partitioning structures (KDTree and BallTree) are used to index the training samples.
\item \textbf{Advantages}:
Na\"ive Bayes models train instantly ($O(nd)$) and have exceptionally fast prediction times ($O(d)$). KNN makes no assumptions about the underlying data distribution, has no training cost, and easily models highly non-linear decision boundaries.
\item \textbf{Limitations}:
Na\"ive Bayes assumes complete conditional feature independence, which is rarely true for text features. KNN requires storing the entire training set (high memory complexity) and has high prediction latency ($O(nd)$) when searching for nearest neighbors.
\end{itemize}

\section{Experimental Setup}
\begin{tabular}{ll}
\toprule
Python Version & 3.11.15 \\
Scikit-Learn Version & 1.8.0 \\
IDE & Jupyter Notebook\\
Operating System & macOS (Version 15.0 Darwin) \\
Hardware & Apple Silicon M-Series CPU, 16GB RAM \\
\bottomrule
\end{tabular}

\section{Hyperparameter Selection}
\begin{tabular}{ll}
\toprule
Parameter & Value\\
\midrule
Number of Neighbors ($k$) & 9 \\
Distance Weighting Scheme & distance \\
Distance Metric & manhattan \\
Search Algorithm & auto \\
\bottomrule
\end{tabular}

\section{Results}
\begin{tabular}{lc}
\toprule
Metric & Value\\
\midrule
Accuracy & 0.8838 \\
Precision & 0.8773 \\
Recall & 0.8436 \\
F1-score & 0.8601 \\
ROC-AUC & 0.9507 \\
\bottomrule
\end{tabular}

\subsection*{Performance Comparisons and Complexity Analysis}
To satisfy the experiment requirements, the complete comparative results for all algorithms and parameter searches are detailed in the tables below:

\subsubsection*{Table 1: Na\"ive Bayes Comparison}
\begin{tabular}{lccc}
\toprule
Metric & Gaussian NB & Multinomial NB & Bernoulli NB \\
\midrule
Accuracy & 0.8219 & 0.8719 & 0.8806 \\
Precision & 0.7233 & 0.9503 & 0.9046 \\
Recall & 0.9385 & 0.7359 & 0.8026 \\
F1-score & 0.8170 & 0.8295 & 0.8505 \\
ROC-AUC & 0.8913 & 0.9591 & 0.9569 \\
\bottomrule
\end{tabular}

\subsubsection*{Table 2: KNN Comparison}
\begin{tabular}{ccccc}
\toprule
k & Accuracy & Precision & Recall & F1 \\
\midrule
1 & 0.8817 & 0.8504 & 0.8744 & 0.8622 \\
3 & 0.8806 & 0.8723 & 0.8410 & 0.8564 \\
5 & 0.8838 & 0.8773 & 0.8436 & 0.8601 \\
7 & 0.8849 & 0.8797 & 0.8436 & 0.8613 \\
9 & 0.8773 & 0.8880 & 0.8128 & 0.8487 \\
11 & 0.8827 & 0.8961 & 0.8179 & 0.8552 \\
\bottomrule
\end{tabular}

\subsubsection*{Table 3: Grid Search vs. Randomized Search}
\begin{tabular}{lcc}
\toprule
Parameter & GridSearchCV & RandomizedSearchCV \\
\midrule
Best k & 9 & 7 \\
Metric & manhattan & manhattan \\
Weights & distance & distance \\
Algorithm & auto & auto \\
CV Accuracy & 0.9196 & 0.9158 \\
Execution Time & 7.3040 s & 0.8249 s \\
\bottomrule
\end{tabular}

\subsubsection*{Table 4: KDTree vs. BallTree}
\begin{tabular}{lcc}
\toprule
Metric & KDTree & BallTree \\
\midrule
Accuracy & 0.8838 & 0.8838 \\
Training Time & 0.002211 s & 0.002535 s \\
Prediction Time & 0.076892 s & 0.080568 s \\
\bottomrule
\end{tabular}

\subsubsection*{Table 5: 5-Fold Cross Validation}
\begin{tabular}{ccc}
\toprule
Fold & Na\"ive Bayes (Gaussian) & Best KNN \\
\midrule
1 & 0.8219 & 0.9001 \\
2 & 0.8065 & 0.9250 \\
3 & 0.7946 & 0.9250 \\
4 & 0.8228 & 0.9163 \\
5 & 0.8337 & 0.9217 \\
\midrule
Average & 0.8159 & 0.9176 \\
\bottomrule
\end{tabular}

\subsubsection*{Table 6: Theoretical Time Complexity}
\begin{tabular}{lcc}
\toprule
Algorithm & Training & Prediction \\
\midrule
Na\"ive Bayes & $O(nd)$ & $O(d)$ \\
KNN (Brute) & $O(1)$ & $O(nd)$ \\
KDTree & $O(n \log n)$ & $O(\log n)$ avg \\
BallTree & $O(n \log n)$ & $O(\log n)$ avg \\
\bottomrule
\end{tabular}

\subsubsection*{Table 7: Experimental Time Analysis}
\begin{tabular}{lcc}
\toprule
Algorithm & Training(s) & Prediction(s) \\
\midrule
Gaussian NB & 0.001182 s & 0.000431 s \\
Multinomial NB & 0.000918 s & 0.000305 s \\
Bernoulli NB & 0.001713 s & 0.000616 s \\
Best KNN & 0.000580 s & 0.021613 s \\
\bottomrule
\end{tabular}

\section{Result Visualization}
All evaluation graphs generated during the experiment are included below, along with quantitative inferences.

\begin{figure}[H]
\centering
\includegraphics[width=0.85\linewidth]{5_confusion_matrices.png}
\caption{Confusion Matrix for Each Model}
\label{fig:confusion_matrices}
\end{figure}

\textbf{Inference:}
The confusion matrices show the exact counts of True Positives, True Negatives, False Positives, and False Negatives. Gaussian NB has the highest recall but suffers from high False Positives (111). Multinomial NB has the lowest False Positives (17) but has high False Negatives (103). The Best KNN and Bernoulli NB strike a better balance, showing both relatively high true positives and true negatives with fewer misclassifications.

\begin{figure}[H]
\centering
\includegraphics[width=0.7\linewidth]{6_roc_curves.png}
\caption{ROC Curves Comparison}
\label{fig:roc_curves}
\end{figure}

\textbf{Inference:}
The Receiver Operating Characteristic (ROC) curves measure the trade-off between True Positive Rate and False Positive Rate. Multinomial NB and Bernoulli NB achieve the highest area under the curve (AUC of 0.9591 and 0.9569, respectively), indicating superior discriminative capacity. Best KNN also exhibits a strong AUC of 0.9507, while Gaussian NB lags behind with an AUC of 0.8913.

\begin{figure}[H]
\centering
\includegraphics[width=0.7\linewidth]{7_precision_recall_curves.png}
\caption{Precision-Recall Curves Comparison}
\label{fig:pr_curves}
\end{figure}

\textbf{Inference:}
The Precision-Recall curves capture classifier efficacy under class imbalance. Multinomial NB achieves the highest Average Precision (AP = 0.9582), followed closely by Bernoulli NB (AP = 0.9555) and Best KNN (AP = 0.9329). Gaussian NB has a lower AP of 0.8659. The high AP values confirm that these classifiers maintain high precision even at high recall rates, making them suitable for spam detection.

\begin{figure}[H]
\centering
\includegraphics[width=0.7\linewidth]{8_accuracy_vs_k.png}
\caption{KNN Test Accuracy vs. Number of Neighbors $k$}
\label{fig:accuracy_vs_k}
\end{figure}

\textbf{Inference:}
The plot illustrates  test accuracy as $k$ varies from 1 to 11. The peak accuracy is achieved at $k=7$ (0.8849) on the test set, with $k=5$ close behind. As $k$ increases past 7, there is a general downward trend in accuracy (dropping to 0.8773 at $k=9$ before recovering slightly at $k=11$). This indicates that smaller local neighborhoods capture the boundary features better, while overly large $k$ values introduce noise from distant classes.

\begin{figure}[H]
\centering
\includegraphics[width=0.8\linewidth]{9_cv_accuracy.png}
\caption{5-Fold Cross-Validation Accuracy Comparison}
\label{fig:cv_accuracy}
\end{figure}

\textbf{Inference:}
The cross-validation accuracy chart shows performance across each fold. The Best KNN consistently outperforms Gaussian Na\"ive Bayes across all 5 folds by a margin of approximately 8--12\%. The average CV accuracy for Best KNN is 91.76\% compared to 81.59\% for Gaussian NB, showing that KNN has significantly higher stability and generalization capabilities on this dataset.

\begin{figure}[H]
\centering
\includegraphics[width=0.7\linewidth]{10_training_time_comparison.png}
\caption{Model Training Time Comparison (seconds)}
\label{fig:training_time}
\end{figure}

\textbf{Inference:}
The training time comparison chart shows that all models train in less than 2 milliseconds. Best KNN shows the fastest training time (0.000580 s) because it is a lazy learner and only stores the dataset. The Na\"ive Bayes models require computing class probabilities and conditional probabilities, resulting in slightly higher but still negligible training times.

\begin{figure}[H]
\centering
\includegraphics[width=0.7\linewidth]{11_prediction_time_comparison.png}
\caption{Model Prediction Time Comparison (seconds)}
\label{fig:prediction_time}
\end{figure}

\textbf{Inference:}
In contrast to training time, the prediction time comparison highlights KNN's key drawback. Best KNN takes 0.021613  seconds to make predictions on the test set, which is 35 to 70 times slower than the Na\"ive Bayes variants (which predict in under 0.0006 seconds). This is because KNN must compute distances to all training samples at test time, while Na\"ive Bayes only performs a simple probabilistic product evaluation.

\begin{figure}[H]
\centering
\includegraphics[width=0.85\linewidth]{12_classifier_comparison_barchart.png}
\caption{Performance Metrics Comparison Across Classifiers}
\label{fig:classifier_comparison}
\end{figure}

\textbf{Inference:}
The bar chart groups accuracy, precision, recall, and F1-score across all models. Best KNN provides the most balanced metrics with high accuracy (over 90\% on CV). Multinomial NB offers the highest precision (95.03\%) but suffers from low recall (73.59\%). Gaussian NB exhibits the opposite behavior, offering excellent recall (93.85\%) but poor precision (72.33\%). Bernoulli NB serves as a middle ground with high recall and precision.

\begin{figure}[H]
\centering
\includegraphics[width=0.75\linewidth]{13_gridsearch_heatmap.png}
\caption{GridSearchCV Mean Accuracy Heatmap ($k$ vs. Distance Metric)}
\label{fig:gridsearch_heatmap}
\end{figure}

\textbf{Inference:}
The heatmap visualizes the grid search optimization space across $k$ and distance metrics. The Manhattan distance metric consistently yields higher mean cross-validation accuracies compared to Euclidean and Minkowski metrics for all values of $k$. The maximum cross-validation score is achieved at $k=9$ using Manhattan distance, illustrating the importance of systematic hyperparameter tuning.

\begin{figure}[H]
\centering
\includegraphics[width=0.7\linewidth]{14_randomizedsearch_score_distribution.png}
\caption{RandomizedSearchCV Validation Score Distribution}
\label{fig:random_search_dist}
\end{figure}

\textbf{Inference:}
The histogram displays the  distribution of mean CV accuracy scores for the 15 parameter configurations randomly sampled during Randomized Search. The distribution is left-skewed, showing that the majority of randomly selected configurations result in a strong validation accuracy of over 91\%. The best score achieved is 0.9158 (with $k=7$), which is extremely close to the full grid search optimum of 0.9196.

\section{Discussion}
Discuss observations, strengths, limitations and improvements.
\begin{itemize}
\item \textbf{Observations}: Among the Na\"ive Bayes models, Bernoulli NB achieved the best accuracy (88.06\%) and F1-score (0.8505) because it naturally models features as binary occurrences. Spambase classification is highly sensitive to the presence of specific keywords (presence vs absence), which is well aligned with the Bernoulli model. Multinomial NB achieved the highest precision (95.03\%), producing only 17 False Positives. In contrast, Gaussian NB had a high recall (93.85\%) but poor precision (72.33\%) due to its normal distribution assumption which conflicts with sparse, right-skewed frequencies. The optimized KNN model  ($k=9$, Manhattan distance, distance-weighted)  achieved the best generalization with a cross-validation accuracy of 91.96\%.
\item \textbf{Strengths}: Grid search effectively located parameters ($k=9$, Manhattan metric, distance weights) that maximize prediction consistency. Using spatial trees (KDTree/BallTree) decreases prediction latency for low-to-medium dimensional spaces.
\item \textbf{Limitations}: KNN's high prediction latency (0.0216 s vs. 0.0004 s for Gaussian NB) represents its main limitation due to the $O(nd)$ calculation. Also, Euclidean distance degrades in higher dimensions ($d=57$), making Manhattan distance ($L_1$ norm) more robust.
\item \textbf{Improvements}: Distance-weighted KNN voting handles higher $k$ values more gracefully. For large scale text data, Na\"ive Bayes is computationally superior due to its $O(d)$ prediction time.
\end{itemize}

\section{Conclusion}
We successfully implemented and compared Na\"ive Bayes and KNN on the Spambase dataset. The optimized KNN model ($k=9$, Manhattan metric, distance weighting) achieved the highest cross-validation accuracy of 91.96\%. Among Na\"ive Bayes models, Bernoulli NB performed best with 88.06\% accuracy. The results confirm a distinct tradeoff: KNN yields higher classification accuracy but requires substantial inference time and storage complexity, whereas Na\"ive Bayes has negligible inference latency and memory requirements at the cost of a slightly lower accuracy.

\section{References}
\begin{enumerate}[label={[\arabic*]}]
\item F. Pedregosa et al., ``Scikit-learn: Machine learning in Python,'' \textit{Journal of Machine Learning Research}, vol. 12, pp. 2825-2830, 2011.
\item A. G\'eron, \textit{Hands-On Machine Learning with Scikit-Learn, Keras, and TensorFlow}, 3rd ed. Sebastopol, CA: O'Reilly Media, 2022.
\item C. M. Bishop, \textit{Pattern Recognition and Machine Learning}. New York, NY: Springer, 2006.
\item UCI Machine Learning Repository, ``Spambase Dataset,'' [Online]. Available: \url{https://archive.ics.uci.edu/ml/datasets/spambase}.
\end{enumerate}

\end{document}
